% 2. Użytkownik i cel
\section{Użytkownik i cel}\label{sec:uzytkownik}

\begin{instrukcja}[wymagane]
Kto korzysta z aplikacji, w jakiej sytuacji i po co (Wykład 1). Opiszcie głównego użytkownika jako
personę, jego cel i dwa-trzy najważniejsze zadania. Odwołajcie się do pojęć z wykładu: cel użytkownika,
skala od sporadycznego do zaawansowanego, wzorce zachowań (np. Satisficing, Pragnienie Natychmiastowej
Satysfakcji).
\end{instrukcja}

\subsection{Persona}
\begin{ramka}{Kasia, 21 lat, studentka III roku informatyki}
\begin{tabularx}{\linewidth}{@{}>{\bfseries}p{3.6cm}L@{}}
Sytuacja & między zajęciami na telefonie sprawdza, co ją czeka w tym tygodniu; wieczorem przy komputerze dopisuje nowe terminy podane przez prowadzących \\
Główny cel & nie przegapić żadnego zaliczenia i wiedzieć, ile czasu zostało na naukę i jaki materiał obowiązuje \\
Doświadczenie & środek skali: korzysta kilka razy w tygodniu, zna aplikację, ale nie chce uczyć się skrótów ani ustawień \\
Frustracje & terminy są rozproszone (Teams, USOS, notatki), a ogólny kalendarz miesza zaliczenia z resztą spraw i nie pokazuje zakresu materiału \\
\end{tabularx}
\end{ramka}

\subsection{Główne zadania użytkownika}
\begin{enumerate}
  \item sprawdzić najbliższy termin i liczbę dni do niego (ekran startowy, 0 kliknięć)
  \item przejrzeć terminy jednego przedmiotu i zakres materiału (lista z filtrem, szczegóły)
  \item dopisać nowy termin podany na zajęciach (formularz „Dodaj termin”)
\end{enumerate}

\subsection{Wzorce zachowań, które uwzględniliśmy}
\textbf{Pragnienie Natychmiastowej Satysfakcji:} najważniejsza informacja (najbliższy termin i liczba
dni) jest na ekranie startowym bez żadnego kliknięcia.
\textbf{Satisficing:} użytkownik klika pierwszą rzecz, która wygląda na dobrą, dlatego cała karta
terminu jest klikalna, a polecenie główne „Dodaj termin” jest jedno i wyraźne.
\textbf{Bezpieczna Eksploracja:} stałe menu i okruszki pozwalają wrócić w każdej chwili, więc można
spokojnie zaglądać w szczegóły.
